%% !TEX program = xelatex
\documentclass[12pt,final]{novelette} % font size, twoside, twocolumn, draft-final, use extarticle for 8-20pt

\title{Diff Eq Final 2544}
\author{Kaka Cring}
\date{\today}

\begin{document}

\makeheading{ตัวอย่างข้อสอบ 2301312 สมการเชิงอนุพันธ์ ปรับปรุงจากข้อสอบไล่ภาคปลาย ปีการศึกษา 2544}

\subsubsection*{PART A}

\begin{question}
    \item {}
    จงหาผลเฉลยในรูปอนุกรมกำลังรอบจุด $x=0$ ของสมการ $2x^2y'' + x(2x+1)y' - y = 0$ โดยวิธีของโฟรเบนิอุส \qmark{12}

    % ----------------------------------------

    % ----------------------------------------

    \item {}
    จงหาค่าของ
    \begin{question}
        \item
        $\displaystyle \int_0^\infty e^{-x^4} \dd{x} \quad \text{(ให้ $t = x^4$)}$ \qmark{3}

        % ----------------------------------------

        % ----------------------------------------

        \item
        $\displaystyle \int_{2\pi}^\infty e^{-4t} \sin 2t \cos 2t \dd{t}$ \qmark{4}

        % ----------------------------------------

        % ----------------------------------------

    \end{question}

    \item {}
    จงเขียนสมการเชิงอนุพันธ์ที่มี $P_4(x)$ เป็นผลเฉลย พร้อมทั้งหา $P_4(x)$ \qmark{3}

    % ----------------------------------------

    % ----------------------------------------

    \item {}
    จงหาผลเฉลยของสมการ $\displaystyle xy'' + y' - \left(\frac1x - \frac x4\right)y = 0$ \qmark{8}

    % ----------------------------------------

    % ----------------------------------------

\end{question}

\subsubsection*{PART B}

\begin{question}[resume]
    \item {}
    จงเติมเฉพาะคำตอบลงในช่องว่างที่กำหนดให้โดยไม่ต้องแสดงวิธีทำ \qmark{5}
    \begin{multicols}{2}
        \begin{question}
            \everymath={\displaystyle}
            \item
            $\mathcal{L}\left\{\frac12\right\} = \ \dotfill$

            \item
            $\mathcal{L}\left\{\sin 2t\right\} = \ \dotfill$

            \item
            $\mathcal{L}\left\{\cos 3t\right\} = \ \dotfill$

            \item
            $\mathcal{L}\left\{e^{-4t}\right\} = \ \dotfill$

            \item
            $\mathcal{L}\left\{t^2\right\} = \ \dotfill$

            \item
            $\mathcal{L}^{-1}\left\{\frac3s\right\} = \ \dotfill$

            \item
            $\mathcal{L}^{-1}\left\{\frac{s}{s^2+9}\right\} = \ \dotfill$

            \item
            $\mathcal{L}^{-1}\left\{\frac{1}{s-5}\right\} = \ \dotfill$

            \item
            $\mathcal{L}^{-1}\left\{\frac{1}{s^3}\right\} = \ \dotfill$

            \item
            $\mathcal{L}^{-1}\left\{\frac{1}{s^2+5}\right\} = \ \dotfill$
        \end{question}
    \end{multicols}

    \pagebreak

    \item {}
    จงทำให้เป็นผลสำเร็จ \qmark{10}
    \begin{multicols}{2}
        \begin{question}
            \everymath={\displaystyle}
            \item
            $\int_0^\infty e^{(4-s)t} \dd{t} = \ ?$

            \item
            $\int_0^\infty e^{-st} (t^2 - 2) \dd{t} = \ ?$

            \item
            $\mathcal{L}\left\{4t^{3/2}\right\} = \ ?$

            \item
            $\mathcal{L}\left\{e^t \cos 2t\right\} = \ ?$

            \item
            $\mathcal{L}\left\{t \sin 3t\right\} = \ ?$
        \end{question}
    \end{multicols}

    \item {}
    จงใช้ทฤษฎีบทสังวัตนาการหา $\displaystyle \mathcal{L}^{-1}\left\{\frac{s^2}{(s^2+1)^2}\right\}$ \qmark{5}

    % ----------------------------------------

    % ----------------------------------------

    \item {}
    จงใช้การแปลงลาปลาซหาผลเฉลยของสมการ $\displaystyle \dv{y}{t} + 2y + 2 \int_0^t y(t-x) \dd{x} = 2t; \quad y(0) = 0$ \qmark{7}

    % ----------------------------------------

    % ----------------------------------------
\end{question}

\subsubsection*{PART C}

\begin{question}[resume]
    \item {}
    จงทำให้เป็นผลสำเร็จ \qmark{18}
    \begin{multicols}{2}
        \begin{question}
            \everymath={\displaystyle}
            \item
            $\mathcal{L}\left\{\int_0^t e^x \cos 2x \dd{x}\right\} = \ ?$

            \item
            $\mathcal{L}\left\{\frac{e^t - e^{-2t}}{t}\right\} = \ ?$

            \item
            $\mathcal{L}^{-1}\left\{\frac{2}{s-3} + \frac{s}{9s^2+1}\right\} = \ ?$

            \item
            $\mathcal{L}^{-1}\left\{\frac{e^{-3s}}{s+1}\right\} = \ ?$

            \item
            $\mathcal{L}^{-1}\left\{\ln \frac{s-1}{s}\right\} = \ ?$

            \item
            $\mathcal{L}^{-1}\left\{\frac{s}{s^2+2s+10}\right\} = \ ?$
        \end{question}
    \end{multicols}

    \item {}
    จงหาผลเฉลยของสมการ $\displaystyle y'' + y =
        \begin{cases}
            t,   & 0 < t < \pi \\
            \pi, & t > \pi
        \end{cases} ; \quad y(0) = 2, \quad y'(0) = 3$ \qmark{7}

    % ----------------------------------------

    % ----------------------------------------

    \item {}
    จงหาผลเฉลยของระบบสมการเชิงอนุพันธ์\\[.5em]
    \hspace*{.5cm}$
        \left\{\begin{aligned}
            (D-1)x + y  & = 2 \sin t         \\[.2em]
            -x + (D+1)y & = 4\cos t - \sin t
        \end{aligned}\right.
    $ \quad เมื่อ $x(0) = 2, \quad y(0) = 5$ \qmark{8}

    % ----------------------------------------

    % ----------------------------------------

\end{question}

\subsubsection*{PART D}

\begin{question}[resume]
    \item {}
    กำหนดฟังก์ชัน
    $f(x) = \begin{cases}
            0,       & -\pi < x < 0 \\
            \pi - x, & 0 < x < \pi
        \end{cases}$ และ $f(x + 2\pi) = f(x)$
    \begin{question}
        \item
        จงเขียนกราฟของฟังก์ชัน $f(x)$ บนช่วง $(-3\pi, 3\pi)$

        % ----------------------------------------

        % ----------------------------------------

        \item
        จงหาอนุกรมฟูเรียร์ของฟังก์ชัน $f(x)$

        % ----------------------------------------

        % ----------------------------------------

        \item
        จงหาผลบวกของ $\displaystyle \sum_{n=1}^{\infty} \frac{1}{(2n - 1)^2}$ \qmark{10}

        % ----------------------------------------

        % ----------------------------------------
    \end{question}

    \item {}
    เส้นลวดมีความยืดหยุ่นเส้นหนึ่งยาว $\pi$ หน่วย ปลายทั้งสองถูกยึดติดแน่นที่จุด $x = 0$ และ $x = \pi$ \\
    เส้นลวดนี้ถูกทำให้เป็นเส้นโค้ง $u(x, 0) = x(\pi - x)$ แล้วปล่อยให้เกิดการสั่นด้วยความเร็ว $u_t(x, 0) = 3$ \\
    จงเขียนปัญหาค่าเริ่มต้นและค่าขอบของโจทย์นี้โดยไม่ต้องหาผลเฉลย (กำหนดให้ $c^2 = 1$) \qmark{3}

    % ----------------------------------------

    % ----------------------------------------

    \item {}
    จงแสดงวิธีหาผลเฉลยของปัญหาค่าเริ่มต้นและค่าขอบซึ่งถูกควบคุมด้วยสมการความร้อนใน $1$ มิติ
    $$
        \pdv{u}{t} = c^2 \pdv[2]{u}{x}, \quad 0 < x < L, \quad t > 0
    $$
    ภายใต้เงื่อนไขขอบ \quad $u(0, t) = 0, \quad u(L, t) = 0, \quad t \geq 0$ \\
    และเงื่อนไขเริ่มต้น \quad $u(x, 0) = f(x), \quad 0 < x < L$ \qmark{12}

    % ----------------------------------------

    % ----------------------------------------

    \item {}
    จงหาผลเฉลยของปัญหาในข้อ 14 เมื่อกำหนดให้ $c^2 = 9, L = 4, f(x) = x + 12$ \qmark{5}

    % ----------------------------------------

    % ----------------------------------------
\end{question}

\subsubsection*{PART E}

\begin{question}[resume]
    \item {}
    จงหาผลเฉลยในรูปอนุกรมกำลังรอบจุดสามัญ $x = 0$ ของสมการ $y'' - xy' - y = 0$ \\
    (เขียนผลเฉลยให้เห็นอย่างน้อยถึงพจน์ $x^5$) \qmark{10}

    % ----------------------------------------

    % ----------------------------------------

    \item {}
    ลวดสปริงเส้นหนึ่งถูกยึดปลายบนติดกับเพดาน
    เมื่อนำวัตถุหนัก $6$ ปอนด์ มาผูกติดไว้ที่ปลายล่างของลวดปริงจะทำให้ลวดสปริงยืดออก $6$ นิ้ว
    ถ้าดึงวัตถุนี้ลงมาต่ำกว่าตำแหน่งสมดุล $4$ นิ้ว แล้วปล่อยด้วยการผลักวัตถุขึ้นด้วยความเร็ว $\frac83$ ฟุต/วินาที \\
    จงหาสมการของการเคลื่อนที่ของวัตถุ พร้อมทั้งบอกค่าแอมพลิจูด และมุมทิ้งช่วง \qmark{10}

    % ----------------------------------------

    % ----------------------------------------

    \item {}
    สำหรับปัญหาการเคลื่อนที่ของวัตถุที่ผูกติดกับปลายลวดสปริงปัญหาหนึ่ง
    ถ้าเราทราบว่าสมการของการเคลื่อนที่ของวัตถุคือ
    $$
        x(t) = \frac13 \sin\left(\pi t - \frac\pi4\right)
    $$
    \begin{question}
        \item
        จงหา แอมพลิจูด คาบ และความถี่ของการเคลื่อนที่ \qmark{1.5}

        % ----------------------------------------

        % ----------------------------------------

        \item
        จงหาตำแหน่ง ความเร็ว ความเร่งของการเคลื่อนที่ เมื่อเวลา $1$ วินาที หลังจากปล่อยให้วัตถุเกิดการเคลื่อนที่
        พร้อมทั้งอธิบายลักษณะของการเคลื่อนที่ของวัตถุ ณ ขณะนั้น \qmark{3.5}

        % ----------------------------------------

        % ----------------------------------------
    \end{question}

    \item {}
    จงหาประจุไฟฟ้า $q(t)$ บนตัวเก็บประจุในวงจรไฟฟ้าแบบอนุกรม $RLC$ ที่เวลา $t = 0.01$ วินาที \\
    เมื่อ $L = 0.05$ เฮนรี, $C = 0.01$ ฟารัด, $E(t) = 0$ โวลต์, $q(0) = 5$ คูลอมบ์ และ $i(0) = 0$ แอมแปร์
    และจงพิจารณาหาเวลาที่ประจุไฟฟ้าบนตัวเก็บประจุเป็นศูนย์ครั้งแรก \qmark{10}

    % ----------------------------------------

    % ----------------------------------------
\end{question}

\end{document}